\chapter{Ánh Trăng Lừa Dối Của Bố Mẹ}
\markboth{Ánh Trăng Lừa Dối Của Bố Mẹ}{Ánh Trăng Lừa Dối Của Bố Mẹ}

\begin{center}
	\textit{\color{bookprimary}``Nhiều phụ huynh không thực sự tin con mình hoàn hảo. Họ chỉ sợ nếu thừa nhận lỗi của con, họ sẽ phải đối diện cả phần bất lực của chính mình.''}
\end{center}

\ngancach

\begin{tcolorbox}[
	colback=bookprimary!6,
	colframe=bookprimary!35,
	arc=5pt,
	title={\bfseries Ghi Chú Mở Chương},
	fonttitle=\bfseries\color{bookprimary}
]
Cha mẹ bênh con là phản xạ rất người. Nhưng phản xạ ấy, nếu thiếu tỉnh táo, sẽ biến thành một dạng ánh sáng lừa dối: soi con theo hướng đẹp nhất, và che mất đúng phần con cần được sửa. Chương này không kết tội phụ huynh; nó mổ xẻ nỗi sợ, sĩ diện, và tình thương mù quáng trong những cuộc đối thoại học đường.
\end{tcolorbox}

\section{Con Tôi Ở Nhà Ngoan Lắm}
\begin{truyen}{Con Tôi Ở Nhà Ngoan Lắm}{Parents in Denial}
\chuhoa{T}{rong buổi gặp riêng, một phụ huynh nói rất chắc: ``Cháu ở nhà ngoan lắm thầy. Tôi nghĩ chắc tại bạn bè trong lớp lôi kéo nên cháu mới quậy vậy thôi.''}

\teacherpair{Dạ, bạn bè có ảnh hưởng thật. Nhưng trong nhóm đó, cháu nhà mình thường là người nghĩ ra trò trước. Nếu mình chỉ đổ cho môi trường mà không nhìn vào phần tính nết và thói quen của con, thì mình đang thương con theo cách làm con khó sửa nhất.}{The teacher says: ``Yes, friends do influence a child. But in that group, your child is often the one starting the game. If we blame only the environment and avoid looking at the child's own habits, we are loving the child in the way that makes change hardest.''}

Chị phụ huynh hơi khựng lại: ``Ở nhà nó hiền lắm mà.''

\teacherpair{Nhiều đứa trẻ ở nhà là một phiên bản, ở lớp là một phiên bản khác. Điều đáng sợ không phải là con mình chưa tốt. Điều đáng sợ là người lớn không chịu tin khi sự thật ấy được nói ra.}{``Many children are one version at home and another at school. What is dangerous is not that your child is imperfect. What is dangerous is when adults refuse to believe the truth once it is spoken.''}

Người mẹ im một lúc rồi thở ra: ``Tại tôi sợ lắm thầy. Mình nghe người khác chê con mình là thấy như họ đang chê mình dạy không ra gì.''

\teacherpair{Em hiểu. Cha mẹ nào cũng đau khi nghe điều xấu về con. Nhưng nếu mình chỉ lo giữ thể diện làm cha mẹ, mình sẽ bỏ lỡ cơ hội sửa cho con lúc còn sớm. Thương con thật không phải là chứng minh nó vô tội. Thương con thật là dám nhìn đúng chỗ nó đang lệch.}{``I understand. Every parent feels pain hearing something bad about their child. But if we focus only on protecting our image as parents, we may miss the chance to correct the child early. Real love is not proving your child innocent. Real love is seeing clearly where they are going off course.''}

\textit{(The parent defends the child using the familiar line of innocence at home. The teacher shifts the focus from image to actual behavior.)}
\end{truyen}

\section{Nó Chỉ Hư Trên Trường Thôi}
\begin{truyen}{Nó Chỉ Hư Trên Trường Thôi}{Splitting the Child in Two}
\chuhoa{C}{ó phụ huynh quả quyết}: ``Ở nhà cháu không vậy đâu. Chắc tại vô lớp gặp mấy bạn kia nên mới thành ra thế.''

Thầy chủ nhiệm đáp rất chậm: ``Dạ, môi trường có ảnh hưởng. Nhưng môi trường không thể rút từ trong đứa trẻ ra một thứ hoàn toàn không có. Nếu ở lớp cháu thích gây chú ý, thích thắng miệng, thích hơn thua, thì ở nhà mình cũng nên nhìn xem những mầm đó đang lớn lên từ đâu.''

Người cha hơi khó chịu: ``Ý thầy là do gia đình tôi?''

\teacherpair{Không ai nói một nguyên nhân duy nhất. Chỉ là giáo dục nguy hiểm nhất khi ai cũng muốn mình vô can. Nhà trường không vô can. Gia đình cũng không vô can. Bản thân đứa trẻ lại càng không vô can.}{``No one is pointing to a single cause. Education becomes most dangerous when everyone wants to be innocent. The school is not innocent. The family is not innocent. The child is not innocent either.''}
\end{truyen}

\section{Bênh Con Trước Mặt Con}
\begin{truyen}{Bênh Con Trước Mặt Con}{Public Defense, Private Harm}
\chuhoa{M}{ột buổi làm việc ba bên}, phụ huynh ngồi cạnh con và nói liên tục: ``Con tôi hiền, nó không hỗn vậy đâu. Chắc cô hiểu lầm.''

Đứa trẻ ngồi im, nhưng ánh mắt sáng lên một kiểu rất lạ: kiểu của người vừa thấy một người lớn đứng chắn toàn bộ hậu quả thay mình.

Cô giáo chủ nhiệm nói nhẹ mà nặng: ``Nếu anh chị muốn trao đổi riêng, em sẵn sàng. Nhưng nếu mình phủ hết lỗi cho con ngay trước mặt con, mình đang dạy nó một bài rất nguy hiểm: chỉ cần có người đủ thương mình, mình không cần đối diện sự thật nữa.''
\end{truyen}

\section{Lòng Thương Thiếu Tỉnh Táo}
\begin{truyen}{Lòng Thương Thiếu Tỉnh Táo}{When Love Stops Helping}
\chuhoa{C}{ó những phụ huynh lo cho con đến mức} làm luôn bài, xin luôn lỗi, giải thích luôn mọi thất bại, và đi đến đâu cũng mở đường trước. Họ nghĩ như vậy là thương.

Một giáo viên lâu năm nói rất gắt: ``Thương kiểu đó giống như thấy con yếu chân mà bế luôn nó cả quãng đời. Vấn đề là một đứa trẻ được bế quá lâu sẽ không còn thấy xấu hổ khi không biết đi nữa.''

Rồi thầy dịu lại: ``Con mình ngã thì đau thật. Nhưng để nó không bao giờ ngã, đôi khi là cách chuẩn bị cho nó một cú ngã rất lớn về sau.''
\end{truyen}

\begin{baihoc}
Bênh con không sai. Nhưng bênh tới mức chặn luôn sự thật thì đang vô tình cản con trưởng thành. Điều trẻ em cần không chỉ là người đứng về phía mình, mà là người dám đứng về phía sự thật để cứu tương lai của mình.
\textit{(Protecting your child is not wrong. Protecting them from the truth is. Children need not only someone on their side, but someone brave enough to stand on the side of truth for the sake of their future.)}
\end{baihoc}

\begin{ghinhoanh}
\textbf{Short English to remember:}

Love should not blind truth.

Protection must still leave room for correction.
\end{ghinhoanh}

\begin{tuhocnhanh}
1. Nếu em là phụ huynh, em có dám nghe sự thật khó chịu về con mình không? \textit{(If you were a parent, could you bear hearing an uncomfortable truth about your child?)}

2. Em đã từng được bênh theo cách làm mình hư thêm chưa? \textit{(Have you ever been defended in a way that made you worse?)}

3. Khác nhau giữa bảo vệ con và che lỗi cho con là gì? \textit{(What is the difference between protecting a child and covering up a child's fault?)}
\end{tuhocnhanh}

\begin{tongketchuong}
Cha mẹ nào cũng muốn là nơi con được quay về. Nhưng nếu nơi quay về đó lúc nào cũng nói con vô tội, con rất dễ mang theo một cái tôi méo mó ra ngoài đời. Tình thương trưởng thành không dịu ở bề mặt; nó sáng ở chỗ dám không nói dối vì con.
\end{tongketchuong}

\begin{goigiaovien}
Khi triển khai chương này trong họp phụ huynh, tuyệt đối nên giảm giọng mỉa. Điều phụ huynh cần nghe không phải là họ tệ, mà là họ đang thương con theo cách nào và hậu quả của cách đó ra sao. Nếu phụ huynh thấy bị xúc phạm, họ sẽ khép lại trước khi hiểu.
\end{goigiaovien}

\ngancach